\documentclass[11pt]{article}
\usepackage[margin=1in]{geometry}
\usepackage{amsmath, amssymb, amsthm}
\usepackage{mathtools}
\usepackage{hyperref}

\newtheorem{theorem}{Theorem}
\newtheorem{lemma}[theorem]{Lemma}
\newtheorem{proposition}[theorem]{Proposition}
\theoremstyle{definition}
\newtheorem{definition}[theorem]{Definition}
\newtheorem{example}[theorem]{Example}

\title{Series and Convergence}
\author{math-foundations concept 16}
\date{}

\begin{document}
\maketitle

\section{Series as limits of partial sums}

\subsection*{First, an example}
Before any formalism, watch a series \emph{converge} numerically. Take the
geometric series with ratio $1/2$:
\[
  \sum_{n=0}^{\infty} \left(\tfrac{1}{2}\right)^n
  \;=\; 1 + \tfrac{1}{2} + \tfrac{1}{4} + \tfrac{1}{8} + \tfrac{1}{16} + \cdots
  \;=\; 2.
\]
\emph{Read aloud:} ``the sum from $n$ equals zero to infinity of one-half to
the $n$ equals two.'' The partial sums march visibly toward $2$:
\[
  s_0 = 1,\quad s_1 = 1.5,\quad s_2 = 1.75,\quad s_3 = 1.875,\quad s_4 = 1.9375,\ \ldots
\]
\emph{Read aloud:} ``$s$-zero equals one, $s$-one equals one point five,'' and
so on. Each step halves the gap to the limit.

Contrast this with the \emph{harmonic series} $\sum_{n=1}^{\infty} 1/n
= 1 + 1/2 + 1/3 + 1/4 + \cdots$, whose terms also shrink to $0$, yet whose
partial sums grow without bound: $s_n \sim \ln n \to \infty$. Vanishing terms
are \emph{necessary} but not \emph{sufficient} for convergence. The rest of
this lesson explains exactly which series fall on each side of that line.

Given a real sequence $(a_n)_{n \ge 0}$, the symbol $\sum_{n=0}^{\infty} a_n$
denotes the formal infinite sum. Its value, when defined, is the limit of the
sequence of partial sums.

\begin{definition}[Series and partial sums]
The \emph{$n$-th partial sum} of $(a_n)$ is
\[
  s_n \;=\; \sum_{k=0}^{n} a_k.
\]
\emph{Read aloud:} ``$s$-sub-$n$ equals the sum from $k$ equals zero to $n$
of $a$-sub-$k$.''
The \emph{series} $\sum_{n=0}^{\infty} a_n$ is the symbolic object associated
with $(a_n)$; we say it \emph{converges to $S \in \mathbb{R}$} iff $s_n \to S$
as $n \to \infty$, and write $\sum_{n=0}^{\infty} a_n = S$. Otherwise the
series \emph{diverges}.
\end{definition}

\begin{definition}[Absolute and conditional convergence]
The series $\sum a_n$ \emph{converges absolutely} iff the series of
non-negative terms $\sum |a_n|$ converges. If $\sum a_n$ converges but
$\sum |a_n|$ diverges, the convergence is \emph{conditional}.
\end{definition}

It is a standard fact that absolute convergence implies convergence: the
partial sums of $\sum a_n$ form a Cauchy sequence, dominated by the Cauchy
tails of $\sum |a_n|$. The converse fails, as witnessed by
$\sum (-1)^{n+1}/n = \ln 2$ paired with the divergent harmonic series.

\begin{example}[Geometric series]
For $r \in \mathbb{R}$ with $|r| < 1$, $\sum_{n=0}^{\infty} r^n = \frac{1}{1-r}$.
Indeed $s_n = (1 - r^{n+1})/(1 - r) \to 1/(1-r)$ as $r^{n+1} \to 0$.
This is the workhorse comparator for the ratio test below.
\end{example}

\section{The ratio test}

The ratio test reduces convergence of $\sum a_n$ to comparison with a
geometric series.

\begin{theorem}[Ratio test]
\label{thm:ratio}
Let $(a_n)$ be a sequence of nonzero real numbers and suppose
\[
  L \;=\; \lim_{n \to \infty} \left| \frac{a_{n+1}}{a_n} \right|
\]
\emph{Read aloud:} ``$L$ equals the limit, as $n$ goes to infinity, of the
absolute value of $a_{n+1}$ over $a_n$.''
exists in $[0, +\infty]$. Then:
\begin{enumerate}
  \item if $L < 1$, the series $\sum_{n=0}^{\infty} a_n$ converges absolutely;
  \item if $L > 1$ (including $L = +\infty$), the series diverges;
  \item if $L = 1$, the test is inconclusive.
\end{enumerate}
\end{theorem}

\begin{proof}
We treat the three cases.

\textbf{Case $L < 1$.} Choose $r$ with $L < r < 1$. Since
$|a_{n+1}/a_n| \to L < r$, there exists $N \in \mathbb{N}$ such that
\[
  \left| \frac{a_{n+1}}{a_n} \right| \;<\; r \qquad \text{for all } n \ge N.
\]
\emph{Read aloud:} ``the absolute value of $a_{n+1}$ over $a_n$ is less than
$r$, for all $n$ greater than or equal to $N$.''
Multiplying these inequalities for $n = N, N+1, \dots, N + k - 1$ gives
\[
  |a_{N+k}| \;<\; r^k\,|a_N| \qquad \text{for every } k \ge 0.
\]
\emph{Read aloud:} ``the absolute value of $a_{N+k}$ is less than $r$ to the
$k$ times the absolute value of $a_N$, for every $k$ at least zero.''
Therefore, for all $n \ge N$,
\[
  \sum_{m=N}^{n} |a_m|
  \;\le\;
  |a_N| \sum_{k=0}^{n-N} r^k
  \;\le\;
  \frac{|a_N|}{1 - r}.
\]
\emph{Read aloud:} ``the sum from $m = N$ to $n$ of $|a_m|$ is at most
$|a_N|$ times the sum from $k = 0$ to $n - N$ of $r^k$, which is at most
$|a_N|$ over $1 - r$.''
The partial sums of the non-negative series $\sum |a_m|$ are bounded above,
hence converge (a monotone bounded sequence has a limit). Adding back the
finitely many terms for $m < N$ shows $\sum |a_n| < \infty$, so $\sum a_n$
converges absolutely.

\textbf{Case $L > 1$.} Choose $r$ with $1 < r < L$ (or $r > 1$ arbitrary if
$L = +\infty$). There exists $N$ such that
$|a_{n+1}/a_n| > r$ for all $n \ge N$, hence $|a_{N+k}| > r^k |a_N|$. Since
$r > 1$, $|a_n| \to \infty$, so $a_n \not\to 0$ and the series diverges by
the term test.

\textbf{Case $L = 1$.} The series $\sum 1/n$ has $L = 1$ and diverges, while
$\sum 1/n^2$ has $L = 1$ and converges. No conclusion can be drawn from $L$
alone.
\end{proof}

\begin{example}[Exponential series]
For every $x \in \mathbb{R}$, setting $a_n = x^n / n!$ gives
\[
  \left| \frac{a_{n+1}}{a_n} \right|
  \;=\;
  \frac{|x|}{n+1}
  \;\longrightarrow\;
  0.
\]
\emph{Read aloud:} ``the absolute value of $a_{n+1}$ over $a_n$ equals
absolute-$x$ over $n + 1$, which tends to zero.''
Thus $L = 0 < 1$, and Theorem~\ref{thm:ratio} yields absolute convergence of
$\sum_{n=0}^{\infty} x^n / n! = e^x$ for every real $x$. The same argument
works for $x \in \mathbb{C}$ and, with $|x|$ replaced by $\|A\|$, for any
square matrix $A$.
\end{example}

\section{Power series}

A \emph{power series} centred at $x_0$ is a series of the form
$\sum_{n=0}^{\infty} c_n (x - x_0)^n$. Its set of convergence is an interval
(possibly a single point or all of $\mathbb{R}$) of \emph{radius of
convergence}
\[
  R \;=\; \frac{1}{\limsup_{n \to \infty} |c_n|^{1/n}}
  \qquad \text{(Cauchy--Hadamard).}
\]
\emph{Read aloud:} ``$R$ equals one over the limit superior, as $n$ goes to
infinity, of the $n$-th root of the absolute value of $c_n$ --- the
Cauchy--Hadamard formula.''
Inside $|x - x_0| < R$ the series converges absolutely; outside it diverges;
behaviour at the endpoints requires a separate test.

This perspective is the bridge from analysis to ML: every smooth activation
function, the matrix exponential, and the leading-order expansion of the
KL divergence are all instances of power-series approximation.

\end{document}
